%% ----------------------------------------------------------------
%% Week-11 problem sheet (lectures: biasVariance, overfitting, symmetry).
%% Every number in the answers was checked in numpy (2026-11-02).
%% ----------------------------------------------------------------
\documentclass{sotonExamBoxes}    % use option [answers] to include them
%% ----------------------------------------------------------------
\usepackage{graphicx}
\usepackage{alltt}
\usepackage{bm}
\usepackage{amssymb}
\usepackage{amsfonts}
\usepackage{amsmath}
\thinmuskip=1mu
%% macros.tex -- shared macro definitions for the problem sheets.
%%
%% Every sheet inputs this instead of carrying its own copy. Input it after
%% the \usepackage lines: the definitions below use bm, amsmath and amssymb.
%%
%% Six macros were defined inconsistently across the sheets when this file
%% was made; the choices are marked CONFLICT below.




\usepackage{stmaryrd}

\renewcommand{\ttdefault}{pcr}

\newcommand{\tr}{\textsf{T}}

\newcommand{\e}[1]{{\rm e}^{#1}}

\newcommand{\E}[1]{{\rm e}^{{\displaystyle #1}}}

\newcommand{\logg}[1]{\log\left(#1\right)}

\newcommand{\Prob}[1]{\mathbb{P}\left(#1\right)}

\newcommand{\Step}{\mathop{\mathrm{O}\hspace{-15pt}\rule[8.5pt]{11pt}{1.2pt}\hspace{2pt}}}

\newcommand\diag{\mathop{\mathrm{diag}}}

\newcommand{\grad}{\bm{\nabla}}

\newcommand{\dd}{\mathrm{d}}

\newcommand{\class}{\mathcal{C}}

\newcommand{\data}{\mathcal{D}}

%% CONFLICT: gaussianProcess.tex had {U}{eur}{b}{n} (Euler bold);
%% every other sheet had this. Matrices in that one sheet now look
%% different -- see the report.

\DeclareMathAlphabet{\mat}{OT1}{cmss}{bx}{n}

%% The identity needs to be distinguishable. In cmss a bold I is a plain
%% vertical bar with no crosspieces, so \mat{I} reads as a rule rather than a
%% letter. Euler's bold I keeps its serifs, so the identity alone is set in
%% Euler while everything else stays cmss.
\DeclareMathAlphabet{\matEuler}{U}{eur}{b}{n}
\newcommand{\Id}{\matEuler{I}}

%% CONFLICT and BUG FIX: 16 sheets defined this using \qpartial,
%% which is defined nowhere, so \pd was broken in all of them.
%% Only problem2.tex had the working \partial form, kept here.

\newcommand{\pd}[2]{\ensuremath{\frac{\partial #1}{\partial #2}}\xspace}

\newcommand{\M}[1]{\ensuremath{\boldsymbol{#1}}\xspace}

\newcommand{\bst}{\rule{0mm}{5mm}}

%% CONFLICT: some sheets omitted the \! . Kept the tighter spacing,
%% which is what every sheet that actually uses \av had.

\newcommand{\av}[2][\,]{\mathbb{E}_{#1}\!\left[ {\strut #2} \right]}

\newcommand{\len}[1]{\| #1 \|}

\newcommand{\inner}[2]{\left\langle #1, #2\right\rangle}

\newcommand{\squeeze}{\setlength{\itemsep}{-2pt}}

%% CONFLICT: two spacings, identical output in maths mode.

\newcommand{\bra}[1]{\left( #1 \right)}

\newcommand{\pred}[1]{\left\llbracket { \small #1} \right\rrbracket}

\newcommand{\mln}{\hat{f}_n(\bm{x})}


\newcommand{\Tr}{\mathop{\mathrm{tr}}\,}

%% NOT SHARED: gaussianProcess.tex redefines \det to take an argument
%% (\det{\mat{V}}), while calculus.tex and vectorSpaces.tex use the standard
%% no-argument operator (\det(\mat{X})). The two cannot coexist, so that
%% redefinition stays local to gaussianProcess.tex.

%% NOT SHARED: \ml and \mm mean different things in different sheets.
%% amgm and convexity use a no-argument \ml and \mm = \hat{m}_c(x);
%% ensembleLearning uses a one-argument \ml and \mm = \hat{m}(x).
%% Both forms are in use, so each sheet keeps its own.

\graphicspath{{figures/}}
\usepackage{parskip}
\AddToHook{env/samepage/begin}{%
\AddToHook{para/before}{\nopagebreak}%
\AddToHook{para/after}{\nopagebreak}%
}
\begin{document}
\unitTitle{Bias, Variance and Generalisation Problem Sheet}
\authors{Adam Pr\"ugel-Bennett}
\unitcode{AICE3001/AICE6002}
\semester{2}
\year{2022}
\duration{45}{0}

\maketitle

This sheet covers the lectures on bias and variance, over-fitting and
symmetry.  Question~1 practises the basic ideas, Question~2 works through
shrinkage in detail and Question~3 ridge regression.  Each ends with a
short Python part (\texttt{import numpy as np} is assumed).  Question~4
is harder and optional: it shows that gradient descent regularises, and
builds an invariant function by averaging over a group.

% ***************************** question 1 **************************

\begin{question}{20}
  \begin{qparts}

    \qpart[4]{The same machine is trained on four different data sets.  At
      one input $\bm{x}$, where the true output is $y = 2.5$, the four
      machines predict $2.0$, $2.4$, $1.6$ and $2.0$.  Find the mean
      machine's prediction, the bias $(\hat{f}_m - y)^2$ and the variance
      at this input, and check that they add up to the average squared
      error.
      \begin{answer}
        $\hat{f}_m = (2.0+2.4+1.6+2.0)/4 = 2.0$, so the bias is
        $(2.0-2.5)^2 = 0.25$.  The variance is
        $\frac14(0^2 + 0.4^2 + 0.4^2 + 0^2) = 0.08$.  The average squared
        error is $\frac14(0.25 + 0.01 + 0.81 + 0.25) = 0.33 = 0.25 + 0.08$.
      \end{answer}}{\lines{5}}

    \qpart[4]{For each pair, which machine would you expect to have the
      higher bias, and which the higher variance: (i) a straight line and
      a degree-15 polynomial fitted to 20 noisy points; (ii) $k$-nearest
      neighbours with $k=1$ and with $k=50$?  Explain briefly.
      \begin{answer}
        (i) The straight line has the higher bias (it cannot follow a
        curved function); the degree-15 polynomial has the higher
        variance (it can pass close to every noisy point, so it changes a
        lot from data set to data set).  (ii) $k=1$ copies the label of
        the single nearest point, so it has low bias and high variance;
        $k=50$ averages over many points, so it has low variance but
        higher bias (it smooths over real detail).
      \end{answer}}{\lines{6}}

    \qpart[3]{You have $1000$ labelled examples and use 5-fold cross
      validation.  How many models do you train, how many examples does
      each see, and how many examples is each validated on?  Why must you
      not measure performance on the training data?
      \begin{answer}
        Five models, each trained on 800 examples and validated on the
        remaining 200; every example is used for validation exactly once.
        Training error is optimistic: a flexible machine can fit
        spurious rules in the training data that do not hold on unseen
        data.
      \end{answer}}{\lines{4}}

    \qpart[3]{An ensemble averages $M$ predictors, each with variance
      $s^2 = 1$ at some input and pairwise correlation $\rho = 0.3$.  Using
      $\mathrm{Var}(\bar{f}_M) = \rho s^2 + (1-\rho)s^2/M$, find the
      variance for $M = 10$ and as $M\to\infty$.  Does averaging change the
      bias?
      \begin{answer}
        For $M=10$: $0.3 + 0.7/10 = 0.37$.  As $M\to\infty$ it falls to the
        floor $\rho s^2 = 0.3$.  Averaging does not change the mean, so the
        bias is unchanged.
      \end{answer}}{\lines{4}}

    \qpart[6]{In Python, take 20 equally spaced inputs on $[0,1]$ and
      targets $\sin(2\pi x) + 0.3\,\epsilon$ with $\epsilon\sim\mathcal{N}(0,1)$.
      For polynomials of degree 1, 3 and 9, fit 1000 independent data sets
      (\texttt{np.polyfit}) and estimate the bias (squared error of the
      mean fit, averaged over 100 test points) and the variance.
      \begin{answer}
\begin{alltt}
rng = np.random.default\_rng(0)\\
x = np.linspace(0, 1, 20)\\
xt = np.linspace(0, 1, 100)\\
f = lambda x: np.sin(2 * np.pi * x)\\
for d in [1, 3, 9]:\\
\hspace*{2em}P = []\\
\hspace*{2em}for \_ in range(1000):\\
\hspace*{4em}y = f(x) + 0.3 * rng.standard\_normal(20)\\
\hspace*{4em}P.append(np.polyval(np.polyfit(x, y, d), xt))\\
\hspace*{2em}P = np.array(P)\\
\hspace*{2em}bias2 = ((P.mean(0) - f(xt))**2).mean()\\
\hspace*{2em}print(d, bias2, P.var(0).mean())
\end{alltt}
        Bias and variance are about $0.207$ and $0.009$ for degree 1,
        $0.005$ and $0.017$ for degree 3, and $0.000$ and $0.043$ for
        degree 9.  The bias falls and the variance rises with the degree;
        degree 3 has the smallest total ($0.022$).
      \end{answer}}{\lines{7}}

  \end{qparts}
\end{question}
\newpage
% ***************************** question 2 **************************

\begin{question}{20}
  We estimate the mean $\theta$ of a distribution with variance
  $\sigma^2$ from $m$ independent samples, whose sample mean is
  $\bar{x}$.  Instead of $\bar{x}$ we use the \emph{shrinkage} estimate
  $\hat{\theta} = c\,\bar{x}$ with $0 \leq c \leq 1$.
  \begin{qparts}

    \qpart[4]{Find $\av{\hat{\theta}}$ and $\mathrm{Var}(\hat{\theta})$.
      For which $c$ is $\hat{\theta}$ unbiased?
      \begin{answer}
        $\av{\bar{x}} = \theta$ and $\mathrm{Var}(\bar{x}) = \sigma^2/m$, so
        $\av{\hat{\theta}} = c\,\theta$ and
        $\mathrm{Var}(\hat{\theta}) = c^2\sigma^2/m$.  It is unbiased only
        for $c = 1$ (or if $\theta = 0$).
      \end{answer}}{\lines{4}}

    \qpart[4]{Show that the mean squared error is
      $\av{(\hat{\theta}-\theta)^2} = (1-c)^2\theta^2 + c^2\sigma^2/m$,
      and identify the bias and variance terms.
      \begin{answer}
        Adding and subtracting $\av{\hat{\theta}} = c\,\theta$,
        \begin{align*}
          \av{(\hat{\theta}-\theta)^2} = \av{(\hat{\theta} - c\theta)^2}
          + (c\theta - \theta)^2 = c^2\frac{\sigma^2}{m} + (1-c)^2\theta^2,
        \end{align*}
        since the cross term $2(c\theta-\theta)\av{\hat{\theta} - c\theta}$
        is zero.  $(1-c)^2\theta^2$ is the bias and $c^2\sigma^2/m$ the
        variance.
      \end{answer}}{\lines{6}}

    \qpart[5]{Find the $c^*$ that minimises the mean squared error and the
      minimum error.  Show that it is smaller than $\sigma^2/m$, the error
      of $\bar{x}$, whenever $\theta\neq0$.
      \begin{answer}
        Setting the derivative $-2(1-c)\theta^2 + 2c\sigma^2/m$ to zero,
        \begin{align*}
          c^* = \frac{\theta^2}{\theta^2 + \sigma^2/m}, \qquad
          \text{error} = \frac{\sigma^2}{m}\,\frac{\theta^2}{\theta^2+\sigma^2/m}
          < \frac{\sigma^2}{m}.
        \end{align*}
        (Substituting $c^*$, the bias is
        $\theta^2(\sigma^2/m)^2/(\theta^2+\sigma^2/m)^2$ and the variance
        $(\sigma^2/m)\theta^4/(\theta^2+\sigma^2/m)^2$; they add to the
        expression above.)  The fraction is less than 1, so a biased
        estimate beats the unbiased one.
      \end{answer}}{\lines{7}}

    \qpart[3]{Why can we not use $c^*$ directly?  What happens to $c^*$ as
      $m\to\infty$, and why does that make sense?
      \begin{answer}
        $c^*$ depends on the unknown $\theta$ we are trying to estimate; in
        practice the amount of shrinkage is chosen by validation (or from
        a prior).  As $m\to\infty$, $\sigma^2/m\to0$ and $c^*\to1$: with
        plenty of data the variance is already small, so there is nothing
        to gain from adding bias.
      \end{answer}}{\lines{4}}

    \qpart[4]{Take $\theta = 2$, $\sigma = 4$, $m = 10$.  Compute $c^*$ and
      the errors for $c = 1$ and $c = c^*$.  Check by simulation in
      Python, drawing $100\,000$ samples of size $m$.
      \begin{answer}
        $\sigma^2/m = 1.6$, so $c^* = 4/5.6 \approx 0.714$, with error
        $1.6\times4/5.6 \approx 1.143$ against $1.6$ for $c = 1$.
\begin{alltt}
rng = np.random.default\_rng(0)\\
theta, sigma, m = 2.0, 4.0, 10\\
xbar = rng.normal(theta, sigma, (100000, m)).mean(1)\\
for c in [1.0, 0.714, 0.5]:\\
\hspace*{2em}print(c, np.mean((c * xbar - theta)**2))
\end{alltt}
        This gives $1.596$, $1.137$ and $1.395$ (theory $1.6$, $1.143$ and
        $1.4$): shrinking too far ($c=0.5$) helps less, but still beats
        $c=1$.
      \end{answer}}{\lines{6}}

  \end{qparts}
\end{question}
\newpage
% ***************************** question 3 **************************

\begin{question}{20}
  Data are generated as $\bm{y} = \mat{X}\bm{w}^* + \bm{\epsilon}$, where the
  noise terms $\epsilon_k$ are independent with mean 0 and variance
  $\sigma^2$.  Ridge regression minimises
  $\len{\mat{X}\bm{w}-\bm{y}}^2 + \eta\len{\bm{w}}^2$.
  \begin{qparts}

    \qpart[5]{In one dimension, $y_k = w^* x_k + \epsilon_k$ for
      $k = 1,\ldots,m$.  With $S = \sum_k x_k^2$, show that
      $\hat{w} = \sum_k x_k y_k/(S+\eta)$ and find its mean and variance.
      For $x = (1, 2, 3)$, $w^* = 0.5$ and $\sigma = 1$, find the bias,
      variance and mean squared error for $\eta = 0$ and $\eta = 4$.
      \begin{answer}
        The derivative $2\sum_k x_k(w x_k - y_k) + 2\eta w$ is zero at
        $\hat{w} = \sum_k x_k y_k/(S+\eta)$.  Since
        $\sum_k x_k y_k = S w^* + \sum_k x_k\epsilon_k$,
        \begin{align*}
          \av{\hat{w}} = \frac{S}{S+\eta}\,w^*, \qquad
          \mathrm{Var}(\hat{w}) = \frac{\sigma^2 S}{(S+\eta)^2}.
        \end{align*}
        Here $S = 14$.  For $\eta = 0$: no bias, variance $1/14 \approx
        0.071$.  For $\eta = 4$: bias $(-4\times0.5/18)^2 \approx 0.012$,
        variance $14/324 \approx 0.043$, total $0.056$, which is smaller.
      \end{answer}}{\lines{8}}

    \qpart[5]{In general write the SVD $\mat{X} = \mat{U}\mat{S}\mat{V}^\tr$,
      with columns $\bm{u}_k$ and $\bm{v}_k$ and singular values $s_k$.
      Show that $\mat{X}^\tr\mat{X} + \eta\mat{I} =
      \mat{V}(\mat{S}^2 + \eta\mat{I})\mat{V}^\tr$, and hence that the
      component of $\hat{\bm{w}}$ along $\bm{v}_k$ is
      $\hat{w}_k = s_k\,\bm{u}_k^\tr\bm{y}/(s_k^2+\eta)$.
      \begin{answer}
        $\mat{X}^\tr\mat{X} = \mat{V}\mat{S}\mat{U}^\tr\mat{U}\mat{S}\mat{V}^\tr
        = \mat{V}\mat{S}^2\mat{V}^\tr$ and $\mat{I} = \mat{V}\mat{V}^\tr$
        ($\mat{V}$ is orthogonal), which gives the first result.  Its
        inverse is $\mat{V}(\mat{S}^2+\eta\mat{I})^{-1}\mat{V}^\tr$, so
        \begin{align*}
          \hat{\bm{w}} = \mat{V}(\mat{S}^2+\eta\mat{I})^{-1}\mat{V}^\tr
          \mat{V}\mat{S}\mat{U}^\tr\bm{y}
          = \mat{V}(\mat{S}^2+\eta\mat{I})^{-1}\mat{S}\,\mat{U}^\tr\bm{y}.
        \end{align*}
        Multiplying by $\bm{v}_k^\tr$ picks out row $k$ of the diagonal
        matrix: $\hat{w}_k = s_k\,\bm{u}_k^\tr\bm{y}/(s_k^2+\eta)$.
      \end{answer}}{\lines{8}}

    \qpart[4]{With $w_k^* = \bm{v}_k^\tr\bm{w}^*$, show that
      $\bm{u}_k^\tr\bm{y} = s_k w_k^* + \bm{u}_k^\tr\bm{\epsilon}$, and hence
      that $\av{\hat{w}_k} = s_k^2 w_k^*/(s_k^2+\eta)$ and
      $\mathrm{Var}(\hat{w}_k) = \sigma^2 s_k^2/(s_k^2+\eta)^2$.
      \begin{answer}
        $\bm{u}_k^\tr\mat{X}\bm{w}^* = \sum_j s_j(\bm{u}_k^\tr\bm{u}_j)
        (\bm{v}_j^\tr\bm{w}^*) = s_k w_k^*$, because the $\bm{u}_j$ are
        orthonormal.  The noise $\bm{u}_k^\tr\bm{\epsilon} = \sum_i
        u_{ik}\epsilon_i$ has mean 0 and variance $\sigma^2\sum_i u_{ik}^2
        = \sigma^2$.  Substituting into part (b) gives both results.
      \end{answer}}{\lines{6}}

    \qpart[6]{Take two components with $s = (2, 0.1)$, $w^* = (1, 1)$ (in
      the $\bm{v}$ basis) and $\sigma = 1$, so that
      $\av{\len{\hat{\bm{w}}-\bm{w}^*}^2} = \sum_k
      \frac{\eta^2 {w_k^*}^2 + \sigma^2 s_k^2}{(s_k^2+\eta)^2}$.  Evaluate it
      for $\eta = 0$ and $\eta = 1$.  Show that each term is smallest at
      $\eta = \sigma^2/{w_k^*}^2$, and check your values by simulation in
      Python.
      \begin{answer}
        $\eta = 0$: $1/4 + 1/0.01 = 100.25$, almost all from the small
        singular value.  $\eta = 1$: $\frac{1 + 4}{25} + \frac{1 +
        0.01}{1.01^2} \approx 0.2 + 0.99 = 1.19$.  Differentiating a term
        $(\eta^2 a + b)/(s^2+\eta)^2$ (with $a = {w_k^*}^2$, $b = \sigma^2s^2$)
        gives a numerator proportional to $\eta a s^2 - b = s^2(\eta a -
        \sigma^2)$, which is zero at $\eta = \sigma^2/a$, whatever $s$ is.
        Here every term is smallest at $\eta = 1$.
\begin{alltt}
rng = np.random.default\_rng(0)\\
U, \_ = np.linalg.qr(rng.standard\_normal((20, 2)))\\
V, \_ = np.linalg.qr(rng.standard\_normal((2, 2)))\\
s = np.array([2.0, 0.1])\\
X = U @ np.diag(s) @ V.T\\
w\_star = V @ np.array([1.0, 1.0])\\
for eta in [0.0, 0.1, 1.0]:\\
\hspace*{2em}err = []\\
\hspace*{2em}for \_ in range(20000):\\
\hspace*{4em}y = X @ w\_star + rng.standard\_normal(20)\\
\hspace*{4em}A = X.T @ X + eta * np.eye(2)\\
\hspace*{4em}w = np.linalg.solve(A, X.T @ y)\\
\hspace*{4em}err.append(((w - w\_star)**2).sum())\\
\hspace*{2em}print(eta, np.mean(err))
\end{alltt}
        The simulation gives $99.8$, $1.88$ and $1.19$, against $100.25$,
        $1.89$ and $1.19$ from the formula.
      \end{answer}}{\lines{8}}

  \end{qparts}
\end{question}
\newpage

% ***************************** question 4 **************************

\begin{question}{25}
  This question is harder and optional.
  \begin{qparts}

    \qpart[6]{We minimise $\frac12\len{\mat{X}\bm{w}-\bm{y}}^2$ by gradient
      descent, $\bm{w}\leftarrow\bm{w} - r\,\mat{X}^\tr(\mat{X}\bm{w}-\bm{y})$,
      starting from $\bm{w}=\bm{0}$, where $\mat{X}$ is $m\times p$ with
      $p > m$ and $\mat{X}\mat{X}^\tr$ invertible.  Explain why $\bm{w}$ is
      always of the form $\mat{X}^\tr\bm{a}$.  If it converges to an exact
      fit, show that the limit is
      $\bm{w} = \mat{X}^\tr(\mat{X}\mat{X}^\tr)^{-1}\bm{y}$.
      \begin{answer}
        Each step adds $-r\,\mat{X}^\tr(\mat{X}\bm{w}-\bm{y})$, which has the form
        $\mat{X}^\tr\bm{b}$; starting from $\bm{0}$, a sum of such vectors
        is still $\mat{X}^\tr\bm{a}$.  An exact fit with $\bm{w} =
        \mat{X}^\tr\bm{a}$ needs $\mat{X}\mat{X}^\tr\bm{a} = \bm{y}$, so
        $\bm{a} = (\mat{X}\mat{X}^\tr)^{-1}\bm{y}$.
      \end{answer}}{\lines{6}}

    \qpart[6]{Show that this is the exact fit with the smallest norm.
      (Hint: any other exact fit is $\bm{w}+\bm{z}$ with
      $\mat{X}\bm{z} = \bm{0}$.)  Why does this matter for over-parameterised
      models such as neural networks?
      \begin{answer}
        If $\mat{X}\bm{z} = \bm{0}$ then $\bm{z}^\tr\bm{w} =
        \bm{z}^\tr\mat{X}^\tr\bm{a} = (\mat{X}\bm{z})^\tr\bm{a} = 0$, so by
        Pythagoras $\len{\bm{w}+\bm{z}}^2 = \len{\bm{w}}^2 + \len{\bm{z}}^2 \geq
        \len{\bm{w}}^2$.  With more parameters than data there are infinitely
        many exact fits; gradient descent from a small start picks a
        small-norm one without being told to.  The optimiser is a
        regulariser, which helps explain why heavily over-parameterised
        models can still generalise (double descent).
      \end{answer}}{\lines{7}}

    \qpart[5]{Check part (a) in Python with $m = 10$, $p = 30$, random
      $\mat{X}$ and $\bm{y}$, learning rate $r = 0.01$ and $20\,000$ steps.
      \begin{answer}
\begin{alltt}
rng = np.random.default\_rng(0)\\
m, p = 10, 30\\
X = rng.standard\_normal((m, p))\\
y = rng.standard\_normal(m)\\
w = np.zeros(p)\\
for \_ in range(20000):\\
\hspace*{2em}w -= 0.01 * X.T @ (X @ w - y)\\
w\_mn = X.T @ np.linalg.solve(X @ X.T, y)\\
print(np.abs(w - w\_mn).max(), np.abs(X @ w - y).max())
\end{alltt}
        Both differences are about $10^{-15}$: gradient descent has found
        the minimum-norm fit.  ($r = 0.01$ is below $2/s_{\max}^2 \approx
        0.024$, so the iteration is stable.)
      \end{answer}}{\lines{6}}

    \qpart[8]{The rotations of a square image by $0$, $90$, $180$ and
      $270$ degrees form a group $C_4$.  For any function $f$ of images,
      show that $\bar{f}(\bm{x}) = \frac14\sum_{k=0}^{3} f(R^k\cdot\bm{x})$, where
      $R$ is a rotation by $90$ degrees, is unchanged when $\bm{x}$ is
      rotated.  Check this in Python with \texttt{np.rot90} for
      $f(\bm{x}) = \sum_{ij} x_{ij}\,(4i+j)$ on a random $4\times4$ image.
      \begin{answer}
        $\bar{f}(R\cdot\bm{x}) = \frac14\sum_{k=0}^3 f(R^{k+1}\cdot\bm{x})$.  As
        $k$ runs over $0,\ldots,3$, $k+1$ runs over $1,\ldots,4$, and
        $R^4 = e$, so this is the same four terms in a different order:
        $\bar{f}(R\cdot\bm{x}) = \bar{f}(\bm{x})$.
\begin{alltt}
f = lambda x: (x * np.arange(16).reshape(4, 4)).sum()\\
def fbar(x):\\
\hspace*{2em}return sum(f(np.rot90(x, k)) for k in range(4)) / 4\\
img = rng.standard\_normal((4, 4))\\
print([f(np.rot90(img, k)) for k in range(4)])\\
print([fbar(np.rot90(img, k)) for k in range(4)])
\end{alltt}
        $f$ takes four different values on the four rotations, while
        $\bar{f}$ takes the same value on all of them.  This is test-time
        augmentation: averaging over a group makes any model invariant.
      \end{answer}}{\lines{8}}

  \end{qparts}
\end{question}

\end{document}
